% Generated draft: requires mathematical and source review.
\documentclass[11pt]{article}
\usepackage[margin=25mm]{geometry}
\usepackage{fontspec}
\setmainfont{Latin Modern Roman}
\usepackage{amsmath,amssymb,mathtools,unicode-math}
\setmathfont{Latin Modern Math}
\usepackage{enumitem,needspace,xurl,hyperref}
\setlength{\parindent}{0pt}
\setlength{\parskip}{5pt}
\setlength{\emergencystretch}{3em}
\newcommand{\N}{\mathbb{N}}
\newcommand{\Z}{\mathbb{Z}}
\newcommand{\Q}{\mathbb{Q}}
\newcommand{\R}{\mathbb{R}}
\newcommand{\C}{\mathbb{C}}
\newcommand{\F}{\mathbb{F}}
\newcommand{\A}{\mathbb{A}}
\newcommand{\PP}{\mathbb{P}}
\newcommand{\E}{\mathbb{E}}
\newcommand{\eps}{\varepsilon}
\newcommand{\dd}{\,\mathrm{d}}
\begin{document}
\section*{UnsolvedMath ID: 20002347}
\textbf{Source problem code:} catalogue entry
\subsection{Definitions and mathematical statement}
Davis's Question 10 has two methodological requests: find a native proof of DPRM over \(\Z\) without referring to \(\N\), and prove DPRM abstractly for a class of rings without reference to \(\N\). The source does not define ``native'' or the class of rings. For a computably presented ring \(R\), use DPRM here to mean that every recursively enumerable finite-arity relation on \(R\) has a positive-existential ring definition, equivalently an existential finite system of polynomial equations, in its stated language. Noncommutative rings use ring equations with their multiplication order respected; no unqualified single-polynomial encoding is assumed.

The aid reports the following signed scalar-transfer theorem. Let \(A\) be a unital \(\Z\)-algebra, finite free as a \(\Z\)-module, with a named additive basis
\[
 e_0=1,e_1,\ldots,e_{d-1}.
\]
Its multiplication is given by fixed integral structure constants, and use the computable signed-coordinate presentation \(A\cong\Z^d\). Suppose the scalar subring \(\Z\cdot1\) is positive-existentially definable without extra parameters in this expanded language. Then the aid states
\[
 \operatorname{DPRM}(A)\quad\Longleftrightarrow\quad
 \operatorname{DPRM}(\Z).
\]
``Parameter free'' here allows the named basis constants already in the language. The expansion of \(A\)-equations into signed integral coordinates handles all ring operations, and the scalar predicate restricts the integer witnesses in the reverse interpretation. The transfer does not use a nonnegative-cone definition or DPRM over \(\N\).

If the center is exactly \(Z(A)=\Z\cdot1\), the scalar predicate is the finite system
\[
 \bigwedge_{i<d}xe_i=e_ix.
\]
Thus, in this precise central finite-free setting with named basis, the aid concludes DPRM from the known integer theorem. No assumption that every finite-free commutative extension has center \(\Z\) is made. The transfer is conditional on DPRM\((\Z)\) and does not itself prove that integer theorem natively.

The supplied literature discussion says Pasten's 2022 framework still obtains DPRM\((\Z)\) through positive-existential bi-interpretation with \(\N\), and that his characterization obstructs a stronger reading requiring structural separation from \(\N\). A structural interpretation characterization does not by itself dictate the exposition of every proof. The exact characterization and its hypotheses are not reproduced, so no blanket impossibility of a new \(\Z\)-internal proof is asserted. The original undefined methodological goal is not marked settled by the transfer theorem.
\subsection{Short English statement}
The source asks for a proof method avoiding appeal to the natural numbers. The aid gives a signed-coordinate DPRM equivalence for finite-free integer algebras with a definable scalar copy, including its central special case. This transfers the known integer theorem and does not supply the requested native proof of it.
\subsection{Sources}
\begin{itemize}
\item Local UnsolvedMath registry, permanent ID 20002347, exact source code catalogue entry; canonical \texttt{aim-logic-notes.json}, \texttt{notes[122]}, tag \texttt{question}; AIM Workshop Problem Lists, Logic, \emph{Extensions of Hilbert's tenth problem}, Question 10 (Davis). Created and updated 14 August 2026.
\item Source URL \url{https://aimath.org/WWN/hilberts10th/hilberts10th.pdf}. The aid reports that the only OCR repair is a missing space/newline between the two sentences. The PDF was not accessed.
\item Pasten's 2022 framework and characterization are named without a full title, publication identifier, theorem number, exact hypotheses, or separate URL.
\item Attempt 1: \texttt{partial\_result}, problem status at run \texttt{partially\_solved}, rigor label \texttt{proved}, novelty confidence \texttt{low}; classification \texttt{PARTIAL-PROGRESS}, registry status \texttt{partially\_solved}. Candidate: signed finite-free scalar transfer. These are dated source labels, not a proof of the first methodological request.
\item Ulam AI / UnsolvedMath Contributors; CC BY 4.0. Machine-generated research aid, not peer reviewed; independently verify before citation.
\end{itemize}
\end{document}
